55\section{Introduction}

The morphological classification of resolved and unresolved sources\footnote{Note that some stars can be
resolved, such as Asymptotic Giant Branch stars, while quasars and galaxies can be unresolved.}, commonly
referred to as ``star-galaxy separation'', is a fundamental step in the analysis of wide-field imaging surveys
such as the Rubin's Legacy Survey of Space and Time (LSST). At bright magnitudes this problem is
relatively simple, since stars outnumber galaxies and the galaxies that do exist are obviously extended on the sky.
Star-galaxy separation becomes increasingly harder at faint magnitudes, such as those probed by
Rubin's Legacy Survey of Space and Time (LSST; for details see \citealt{2019ApJ...873..111I} and references therein)
because galaxy counts rapidly increase with magnitude, while their angular sizes become smaller.

A variety of methods have been proposed in the literature for morphological classification of sources in images. 
\cite{2020AJ....159...65S} (hereafter SIL2020) discussed the statistical foundations of morphological star–galaxy separation and showed
that many of the star-galaxy separation metrics in common use today, e.g., by Sloan Digital Sky Survey (SDSS) {\it photo}
pipeline \citep{2002SPIE.4836..350L}, or SExtractor pipeline \citep{1996A&AS..117..393B}, are closely
related both to each other, and to the Bayesian model odds ratio, first derived by \cite{1979AJ.....84.1526S}.
SIL2020 quantified the scaling of various algorithms with the source noise properties and showed that the Bayesian model
odds ratio achieves the best peformance. They also showed how to probabilistically combine multiple measurements,
either of the same type (e.g., subsequent exposures), or differing types (e.g., multiple bandpasses), or differing
methodologies (e.g., morphological and color-based classification).

Following SDSS, LSST photometric catalogs use a binary classifier \texttt{extendedness} in each band, based on the
threshold difference between PSF, \(m_{\mathrm{PSF}}\), and CModel, \(m_{\mathrm{CModel}}\),
magnitudes, \(\Delta m = m_{\mathrm{PSF}} - m_{\mathrm{CModel}}\). For unresolved sources, this \(\Delta m \)
is vanishing within noise, with \texttt{extendedness=0}, while for resolved objects, selected by \(\Delta m >
0.016\) mag, \texttt{extendedness=1}. Resolved objects have positive \(\Delta m \) because their
PSF flux is smaller than CModel flux due to a narrow weighting PSF profile. For an illustration of \(\Delta m \)
distribution in Rubin's Data Preview 1 catalogs, see Fig.~28 in \cite{2026AJ....171..360V}. 

Although  \texttt{extendedness} was successfuly used in thousands of papers based on SDSS catalogs, this
approach is not optimal for LSST. As discussed in SIL2020, when the star-galaxy separation threshold is optimized
at the faint, low S/N, end of the data, the ability to recognize barely resolved objects at the bright end is not fully exploited. 
More importantly, the binary nature of \texttt{extendedness} parameter prevents optimal combination of information from
different bands, as well as from other classification methods (e.g., based on colors, or other datasets). SIL2020 showed how
to turn \(\Delta m \) into a probability distribution, which then allows proper probabilistic combination of multiple bands.
With appropriately chosen priors for star-to-galaxy counts ratio, the posterior probability that a source is unresolved
(or resolved), can be computed in Bayesian framwork. 

The main goal of this paper is to implement the steps towards Bayesian star-galaxy separation outlined in SIL2020
using Rubin's Data Preview 2 data. We describe our dataset and method in \S 2, and quantify its performance using
HST-based labels in \S 3. We discuss our results and how to retrieve our catalog with improved star-galaxy separation in \S 4. 
 
